\documentclass[11pt]{article}
\usepackage{amsmath}
\usepackage{geometry}
\geometry{a4paper, margin=1in}

\title{\textbf{Computational Evaluation of Algebraic and Geometric Series Using Iterative Loops}}
\author{Aparajita Chakraborty \\ Bsc Physics Honors}
\date{27 May, 2026 }

\begin{document}

\maketitle

\section{Introduction}
This report presents the mathematical formulation and algorithmic implementation for evaluating finite series expansions. Specifically, we examine the computational execution of a cumulative summation loop, and a sequential product loop, designed to run natively without reliance on external mathematical packages.

\section{Mathematical Formulations}
In computational physics, continuous mathematical expressions are often evaluated discretely using bounded loops.

\subsection{Summation Series}
A standard finite arithmetic or geometric progression can be represented by the summation of discrete terms where $i$ is in the loop index:
\begin{equation}
	S_n = \sum_{i=1}^{n} f{i}
\end{equation}
The algorithm initializes an accumulator variable $S = 0$ and sequentially adds the value of $f(i)$ at each step.

\subsection{Product Series}
Similarly, a bounded product series (such as a factorial sequence or product expansion) multiplies terms sequentially:
\begin{equation}
	P_n = \prod_{i=1}^{n} f(i)
\end{equation}
The computational tracking variable must be initialized to $P = 1$ ( the multiplicative identity) to avoid annihilating the product array during iteration.

\section{Conclusion}
By utilizing primitive conditional loops, we successfully mapped formal mathematical $\sum$ and $prod$ operations into executable Python logic, providing a foundation for evaluating more complex discrete physical systems.
\end{document}